% Laget automatisk av elevark/lag_elevark.py fra sim/magnetfelt-bane.html. Endringer her blir overskrevet.
\documentclass[11pt]{article}
\usepackage{elevark}
\begin{document}
\elevarkhode{Fysikk 2 · Magnetiske felt og krefter}{Ladd partikkel i magnetfelt}{https://fysikklab.no/sim/magnetfelt-bane.html}{Simuleringen}

En ladd partikkel som beveger seg vinkelrett på et magnetfelt, får en kraft $F = qvB$ som hele tiden står vinkelrett på farten. Da går partikkelen i en sirkel. I et massespektrometer bruker vi dette til å skille atomer som bare har litt ulik masse.

\section*{Oppgaver}

Start med et proton, $U = 500\,\mathrm{V}$ og $B = 50\,\mathrm{mT}$ inn i planet.

\begin{oppgaveliste}
\oppgave \fase{gjett}{Gjett}Protonet går rett opp inn i feltet. Bøyer det mot høyre eller venstre? Bruk høyrehåndsregelen, og test. Hva skjer hvis du snur feltet eller bytter til et elektron?
\svar
\oppgave \fase{test}{Regn}Regn ut farten til protonet, den magnetiske kraften og radien i banen. Sjekk med simuleringen.
\svar
\oppgave \fase{gjett}{Gjett}Hva skjer med radien hvis du gjør spenningen fire ganger så stor? Hva om du i stedet dobler magnetfeltet? Test begge.
\svar
\oppgave \fase{forklar}{Forklar}Den magnetiske kraften endrer ikke farten til partikkelen, bare retningen. Hvorfor? Bruk begrepet arbeid.
\svar
\oppgave \fase{test}{Regn}Vis at omløpstiden er $T = 2πm / (qB)$. Regn ut hvor lang tid protonet er i feltet. Hva skjer med tiden hvis du øker spenningen? Test.
\svar
\oppgave \fase{test}{Undersøk}Bytt til et elektron. Hvorfor blir sirkelen så mye mindre? Hvor svakt må feltet være for at elektronet skal få radius 7,5~cm? Regn først, og test så.
\svar
\oppgave \fase{test}{Regn}Velg massespektrometer og klor, $U = 2000\,\mathrm{V}$ og $B = 200\,\mathrm{mT}$. Klor har to isotoper, ${}^{35}\mathrm{Cl}$ og ${}^{37}\mathrm{Cl}$. Regn ut hvor langt fra spalten hver isotop treffer detektoren, og hvor langt det er mellom treffpunktene.
\svar
\oppgave \fase{utfordring}{Utfordring}Velg uran med $U = 2000\,\mathrm{V}$ og $B = 500\,\mathrm{mT}$. Hvor langt er det mellom treffpunktene for ${}^{235}\mathrm{U}$ og ${}^{238}\mathrm{U}$? Hvorfor er det mye vanskeligere å skille uranisotoper enn klorisotoper? Finn en sammenheng mellom avstanden og massene.
\svar
\end{oppgaveliste}

\end{document}
